\documentclass{article}
\usepackage{graphicx}   % images
\usepackage{amsmath}    % equations
\usepackage{booktabs}   % nicer tables
\usepackage{hyperref}

\title{myoL3: An Automated Pipeline for L3-Level Trunk-Muscle
Segmentation and Body-Composition Quantification from CT}
\author{Yousif Al-Khoury}
\date{August 2026}

\begin{document}

\maketitle

% =====================================================================
\section{Introduction}
% Bullet-point structure only — expand into prose later.
\begin{itemize}
  \item Clinical motivation: skeletal-muscle wasting (sarcopenia / ICU-acquired
        weakness) is common and prognostic; muscle quantity and quality at the
        L3 vertebral level are established surrogates for whole-body composition.
  \item Why L3: single axial slice / short slab at L3 correlates with total
        skeletal-muscle mass; standard in body-composition research.
  \item Current bottleneck: manual or semi-automatic L3 selection and muscle
        tracing is slow, requires expertise, and is a barrier to large cohorts
        and to longitudinal (repeated-scan) analysis.
  \item Gap: need a fully automated, reproducible pipeline that (i) finds L3 in
        an arbitrary CT, (ii) segments the trunk muscles, and (iii) reports both
        muscle \emph{quantity} (area) and \emph{quality} (fat infiltration).
  \item Contribution: \textbf{myoL3}, an installable pipeline that localizes L3,
        segments psoas / quadratus lumborum / erector spinae--multifidus,
        strips intramuscular fat, and outputs per-muscle, per-side
        body-composition metrics; distributed as a package with batch/HPC
        tooling for cohort-scale processing.
  \item State the specific aims / questions addressed by this paper.
\end{itemize}

% =====================================================================
\section{Methods}

\subsection{Study cohort and imaging}
% TEMPLATE — fill cohort specifics.
Describe the source cohort, inclusion/exclusion criteria, number of patients
and scans, timepoints (e.g.\ admission, interim, discharge), scanner(s), and
acquisition parameters. State how reference (ground-truth) L3 bounds and muscle
segmentations were obtained for the validation subset, and by whom.

\subsection{Pipeline overview}
myoL3 processes a single CT volume through four stages: (1) resampling to a
fixed working resolution, (2) L3-level localization and axial cropping, (3)
trunk-muscle segmentation on the crop, and (4) fat stripping and
body-composition quantification. Each stage is described below. The pipeline is
implemented in Python and distributed as an installable package; model weights
are retrieved from a public repository on first use.

\subsection{Preprocessing}
Each input volume is reoriented to a common anatomical convention (RAI, so that
the axial axis is superior--inferior) and resampled to an isotropic in-plane and
near-isotropic through-plane target spacing of
0.5\,mm\,$\times$\,0.5\,mm\,$\times$\,0.625\,mm
using linear interpolation, matching the resolution the models were trained on.
Resampling to a fixed spacing normalizes heterogeneous acquisitions and yields
directly comparable cross-sectional areas across scans.

\subsection{L3 localization}
L3 localization is posed as a one-dimensional boundary-regression problem along
the superior--inferior ($z$) axis. Each axial slice is windowed to a
soft-tissue range and resized in-plane to a fixed size, then encoded by a
Residual Network (ResNet-18) slice encoder. A bidirectional Gated Recurrent Unit
(GRU) models the sequence of slice features along $z$, giving each slice a
context-aware representation. Three heads are attached: two boundary heads that,
via a softmax over $z$ and a soft-argmax (expected-index) operation, predict the
continuous superior and inferior boundaries of the L3 body, and a per-slice
\emph{presence} head that predicts whether each slice belongs to L3. The
soft-argmax construction guarantees a single, ordered, non-empty interval by
design.

Because full CT volumes are longer than the training window, inference uses a
sliding window over $z$: each window's boundary distributions are accumulated
into global-$z$ votes, weighted by the window's presence probability (which
suppresses windows over the chest, pelvis, or legs) and by the sharpness of its
boundary distributions. The final bounds are selected by anchoring on the
presence peak---the most reliable signal---and taking the strongest boundary
votes on each side within a physical span cap of 43.75\,mm
($\approx$L3 body height), which prevents spurious distant votes from stretching
the crop. The volume is then cropped in $z$ to the predicted L3 slab; in-plane
field of view is preserved.

\subsection{Muscle segmentation}
The L3 crop is segmented with a 3D UNet (nnU-Net, full-resolution
configuration) that labels the psoas, quadratus lumborum, and erector
spinae--multifidus. The trained network is loaded from a self-contained
checkpoint and driven in-process (the network's own plans, preprocessing, and
sliding-window tiled inference with Gaussian weighting and test-time mirroring),
so segmentation preprocessing is identical to training. The label map is
returned on the crop grid.

\subsection{Fat stripping and compartmental quantification}
For each muscle, the mask is eroded in-plane by two voxels to define an interior
``core'', separating the true muscle interior from the one-to-two-voxel boundary
rim that reads as fat by partial-volume averaging. Voxels are then classified by
attenuation into four compartments:
\begin{itemize}
  \item \textbf{Muscle} --- core voxels at or above an adaptive
        muscle/low-attenuation threshold (see below).
  \item \textbf{Fat--partial-volume (low-attenuation) muscle} --- core voxels
        between $-30$\,HU and the adaptive threshold.
  \item \textbf{Intramuscular fat} --- core voxels below $-30$\,HU (the
        standard intramuscular-adipose-tissue threshold).
  \item \textbf{Outer-edge partial-volume fat} --- the eroded boundary rim.
\end{itemize}
The muscle/low-attenuation boundary is set \emph{per scan and per muscle} by
fitting a two-component Gaussian mixture model (GMM) to that muscle's core
attenuation histogram (by expectation--maximization) and taking the crossover
between the two Gaussians. This adapts the threshold to each patient's muscle
attenuation, which is shifted downward in myosteatosis; the fixed
$-30$\,HU fat threshold is retained for cross-scan comparability of the fat
compartment. The \emph{total-muscle} segmentation is the original mask with the
intramuscular-fat voxels removed; this is the final segmentation used for
area quantification.

\subsection{Body-composition metrics}
Left and right sides are separated at the patient midline, defined as the
$x$-centroid of all muscle voxels. For each muscle and side, per-slice
cross-sectional area (voxel count $\times$ in-plane pixel area) is computed for
each compartment. For the total-muscle cross-sectional area, contiguous outlier
slices at the superior and inferior edges of the muscle's $z$-extent---where
tapering and partial-volume effects make segmentation unreliable---are trimmed
using a robust (median/median-absolute-deviation) criterion before averaging.
Reported per muscle $\times$ side are the mean, median, and standard deviation of
the total-muscle cross-sectional area (edge-trimmed) and of the three fat
compartments (untrimmed), all in mm$^2$, together with the
adaptive threshold used.

\subsection{Implementation and batch processing}
The pipeline is packaged as an installable Python tool with a command-line
interface and programmatic API; model weights are hosted on a public model
repository and fetched by an installer. For cohort-scale processing, a batch
layer submits one job per scan on a Slurm HPC cluster and a compilation step
aggregates the per-scan metrics into cohort-level tables (a wide muscle-area
table and a tidy long table of all compartments).

\subsection{Statistical analysis}
% TEMPLATE — fill/confirm.
Describe the analyses: descriptive statistics of the cohort; agreement between
automated and reference measurements (validation); longitudinal comparisons
across timepoints (e.g.\ paired tests of muscle area between admission and
discharge); left/right comparisons; and correlation analyses (e.g.\ muscle loss
vs.\ fat change). State the outlier-handling procedure (e.g.\ removal of
segmentation failures and influential points by Cook's distance), the software
used, and the significance threshold.

% =====================================================================
\section{Results}
% Bullet-point structure + validation template.
\begin{itemize}
  \item \textbf{Cohort description.} Number of patients/scans per timepoint;
        demographics; scan-quality summary; number excluded and why.
  \item \textbf{L3 localization validation.} Agreement of predicted vs.\
        reference L3 bounds (e.g.\ boundary error in slices/mm, interval IoU) on
        the held-out set. \emph{[Table/figure placeholder.]}
  \item \textbf{Muscle segmentation validation.} Overlap of automated vs.\
        reference masks per muscle (e.g.\ Dice, Hausdorff). \emph{[Table
        placeholder.]}
  \item \textbf{Cross-sectional-area agreement.} Automated vs.\ reference CSA
        per muscle (Bland--Altman / correlation). \emph{[Figure placeholder.]}
  \item \textbf{Longitudinal muscle change.} Muscle area across timepoints;
        paired admission$\rightarrow$discharge change per muscle and total.
        \emph{[Table placeholder.]}
  \item \textbf{Body-composition / fat.} Fat-compartment findings and their
        relationship (or lack thereof) to muscle change. \emph{[Placeholder.]}
  \item \textbf{Left/right.} Side differences and asymmetry, if reported.
  \item \textbf{Runtime / feasibility.} Per-scan time, cohort throughput.
\end{itemize}

\begin{table}[h]
\centering
\caption{Validation of automated measurements against reference (template).}
\label{tab:validation}
\begin{tabular}{lcccc}
\toprule
Muscle & Dice & Boundary err.\ (mm) & CSA bias (mm$^2$) & CSA LoA \\
\midrule
Psoas               & -- & -- & -- & -- \\
Quadratus lumborum  & -- & -- & -- & -- \\
Erector spinae--multifidus & -- & -- & -- & -- \\
\bottomrule
\end{tabular}
\end{table}

% =====================================================================
\section{Discussion}
% Bullet-point structure only.
\begin{itemize}
  \item \textbf{Summary of findings} relative to the aims.
  \item \textbf{Comparison} with existing automated L3 / body-composition tools;
        what myoL3 adds (end-to-end localization + segmentation + adaptive fat
        quantification; open, installable, HPC-ready).
  \item \textbf{Methodological note.} The per-scan adaptive GMM threshold makes
        the muscle/low-attenuation compartments \emph{not directly comparable}
        across scans with different thresholds; fixed-threshold intramuscular
        fat and mean muscle attenuation are the comparable quality metrics.
  \item \textbf{Limitations.} Sensitivity to input resolution (upsampling coarse
        scans yields interpolated, not real, detail); reliability of the
        smallest muscle (quadratus lumborum); single-institution / retrospective
        design; validation-set size.
  \item \textbf{Future work.} Broader multi-site validation; scan-quality
        screening; extension to additional levels/muscles; prospective
        clinical-outcome association.
  \item \textbf{Conclusion.} One-paragraph takeaway.
\end{itemize}

% =====================================================================
% \section*{Data and code availability}
% myoL3 is available as a package (installation, model weights, and batch
% tooling). State repository/DOI here.

\end{document}
